% Responsable: Naranjo Alfonzo Anthony Alexander
% FORM-PC-01 · Rondas 1, 2 y 3: escenario, elección y riesgos.
\subsection{Aplicación móvil de un emprendimiento}
\label{sec:escenario-app}

\subsubsection{Caracterización}

El enunciado únicamente indica una aplicación móvil de un emprendimiento,
sin precisar qué vende, su presupuesto o las condiciones de desarrollo.
Por ello, se plantea un \textbf{escenario hipotético}: un pequeño
emprendimiento de productos artesanales quiere ofrecer su catálogo y recibir
pedidos desde una aplicación móvil. Los datos siguientes son supuestos de
análisis, no resultados de entrevistas con un negocio real.

\paragraph{Descripción y alcance}
Se asume que la aplicación permitirá a los clientes consultar productos,
ver precios y disponibilidad, preparar una solicitud de pedido y conocer su
estado. El propietario podrá actualizar el catálogo y confirmar o rechazar
pedidos. En la primera versión, el pago y la entrega se coordinarán fuera de
la aplicación. No se incluirán pasarela de pagos, facturación electrónica,
servicio propio de reparto ni integración con sistemas empresariales.

\paragraph{Cliente y actores del proceso}
El cliente es la persona propietaria del emprendimiento, quien establece las
prioridades y acepta las funcionalidades. Los compradores son usuarios
principales y ayudan a evaluar la facilidad de uso. Se asume que una persona
encargada de atender pedidos participa en la definición del flujo de
confirmación; en un negocio pequeño puede ser la misma propietaria. El
equipo de desarrollo obtiene, documenta, negocia y valida los requisitos.

\paragraph{C1. Estabilidad de los requisitos: baja}
Aunque el objetivo de mostrar productos y recibir pedidos está claro, se
supone que la organización del catálogo, el proceso de compra, las
notificaciones y la presentación de las pantallas cambiarán después de que
el propietario y los clientes prueben las primeras versiones.

\paragraph{C2. Criticidad del sistema: baja}
Una falla podría provocar pedidos perdidos, molestias a los clientes o una
venta no concretada, pero el escenario no incluye control de equipos
peligrosos ni procesamiento de pagos dentro de la aplicación. La valoración
baja no elimina la necesidad de proteger los datos de contacto y verificar
que cada pedido quede registrado correctamente.

\paragraph{C3. Disponibilidad del cliente: alta}
Se asume que la persona propietaria puede responder consultas por mensajería
durante la semana y revisar una demostración breve al cerrar cada ciclo. Se
programarán sesiones puntuales con compradores para recibir opiniones sobre
el catálogo y la solicitud de pedidos.

\paragraph{C4. Tamaño y complejidad: pequeño}
El alcance inicial se limita al catálogo, los pedidos, su estado y la
gestión básica por parte del propietario. Participan pocos tipos de usuario
y no se contemplan integraciones externas complejas. Se supone un equipo de
desarrollo reducido.

\paragraph{C5. Regulación y contrato: baja}
Se asume que el proyecto no depende de una contratación pública ni de un
contrato que congele todas las funciones desde el inicio. Las prioridades
podrán acordarse directamente con la persona propietaria. Esto no significa
que no existan obligaciones legales: antes de operar habrá que revisar las
que correspondan al tratamiento de datos personales y a la actividad
comercial del emprendimiento.

\paragraph{C6. Documentación exigida: baja}
Se considera suficiente mantener historias de usuario, criterios de
aceptación, un registro de decisiones, evidencias básicas de prueba y una
guía de uso. No se supone una auditoría contractual que requiera informes
extensos para aceptar cada función.

\begin{table}[H]
\centering\footnotesize
\caption{Caracterización del escenario hipotético de la aplicación móvil}
\label{tab:caracterizacion-app}
\begin{tabularx}{\textwidth}{@{}Yp{0.20\textwidth}@{}}
\toprule
\textbf{Criterio} & \textbf{Valoración}\\
\midrule
C1. Estabilidad de los requisitos & Baja\\
C2. Criticidad del sistema & Baja\\
C3. Disponibilidad del cliente & Alta\\
C4. Tamaño y complejidad & Pequeño\\
C5. Regulación y contrato & Baja\\
C6. Documentación exigida & Baja\\
\bottomrule
\end{tabularx}
\end{table}

\subsubsection{Modelo elegido y justificación}

\textbf{Modelo elegido: ágil, con entregas incrementales y revisiones
frecuentes.} La matriz de la Tabla~\ref{tab:matriz-app} le asigna 18 puntos,
frente a 14 del iterativo y 9 de cascada. La selección depende de que se
cumpla la disponibilidad del propietario planteada como supuesto.

\begin{enumerate}
  \item \textbf{Requisitos cambiantes:} las pantallas y el recorrido para
  hacer pedidos pueden ajustarse cuando los compradores usen prototipos o
  versiones iniciales. El enfoque ágil permite revisar y priorizar esas
  necesidades continuamente \cite{schon2017}.

  \item \textbf{Participación del propietario:} las revisiones breves y las
  respuestas frecuentes permiten decidir rápidamente qué mejorar o qué
  función tiene más valor para el negocio.

  \item \textbf{Alcance pequeño y divisible:} se puede comenzar mostrando
  un catálogo funcional, añadir la solicitud de pedidos y luego incorporar
  la consulta de estados. Cada entrega permite observar problemas reales
  sin esperar a terminar toda la aplicación.

  \item \textbf{Pocas restricciones contractuales y documentales:} el
  emprendimiento puede reajustar prioridades sin tramitar modificaciones
  formales extensas, manteniendo de todas maneras criterios de aceptación,
  pruebas y registro de cambios.
\end{enumerate}

\textbf{Por qué no cascada como modelo principal.} Aunque facilitaría
establecer de antemano una especificación y un plan, exigiría decidir desde
el inicio detalles de experiencia de usuario que todavía no se conocen. Si
las pruebas con clientes muestran que el flujo de pedidos es confuso,
corregirlo tardíamente podría causar retrabajo.

\textbf{Por qué no el iterativo como modelo principal.} También es una
alternativa válida, ya que permite refinar requisitos y evaluar prototipos.
Sin embargo, bajo los supuestos de alta disponibilidad del propietario y
funciones pequeñas que pueden ponerse a prueba, resulta más conveniente
integrar las revisiones de requisitos con la construcción y entrega
frecuente de funcionalidades. El iterativo sería una opción más fuerte si
las revisiones solo pudieran efectuarse en hitos separados o se exigiera
una ERS formal actualizada en cada ciclo \cite{sommerville2016,schon2017}.

\subsubsection{Riesgos del modelo elegido y mitigación}

Los riesgos siguientes corresponden al escenario supuesto; no son problemas
comprobados en un emprendimiento real.

\begin{enumerate}
  \item \textbf{Crecimiento continuo del alcance.} El propietario podría
  solicitar pagos en línea, reparto propio u otras funciones fuera del
  objetivo inicial. \emph{Mitigación:} mantener una lista priorizada de
  historias de usuario, separar el alcance inicial de las mejoras futuras
  y revisar el costo y el tiempo antes de aceptar cambios relevantes.

  \item \textbf{Requisitos confusos o decisiones contradictorias.}
  Al recibir comentarios frecuentes, el equipo podría modificar pantallas
  sin criterios claros. \emph{Mitigación:} definir criterios de aceptación
  para cada historia, registrar decisiones y validar el flujo de pedido con
  el propietario antes de cerrar una entrega.

  \item \textbf{Pérdida de disponibilidad del cliente.} Aunque se supone
  alta participación, las actividades diarias del negocio pueden impedir
  revisar una versión. \emph{Mitigación:} acordar de antemano reuniones
  breves, nombrar a quien pueda responder dudas y no dar por aprobada una
  función que depende de una decisión pendiente.

  \item \textbf{Errores introducidos al agregar funciones.} Una nueva
  pantalla o actualización podría afectar la creación y el seguimiento de
  pedidos. \emph{Mitigación:} realizar pruebas del flujo completo de pedido
  en cada entrega, comprobar el registro de estados y probar la aplicación
  en dispositivos de distintas características.

  \item \textbf{Protección de datos y documentación descuidadas.} La
  rapidez de las entregas podría llevar a almacenar más información de la
  necesaria o a omitir registros de cambios. \emph{Mitigación:} recoger
  solo datos indispensables, establecer permisos de acceso, verificar
  requisitos de privacidad y mantener historias, decisiones y resultados
  de prueba vinculados a cada versión.
\end{enumerate}
